\documentclass[twocolumn]{iopjournal}

\usepackage{amsmath,amssymb,amsfonts,amsthm,graphicx}

\fancyhfoffset[L]{0mm}
\setlength{\marginparwidth}{0pt}
\setlength{\marginparsep}{0pt}

\renewcommand{\articletype}[1]{%
  {\vspace*{-8mm}\noindent \Large \sf Preprint}\par
  \vspace*{6mm}\noindent{\scriptsize \sf{\bfseries \MakeUppercase{#1}}}}

\fancyhead[L]{}
\fancyhead[R]{}

\hypersetup{pdftitle={Where Closure Sits: the Field's Stagnation Surface and the Jacobi Surface}, pdfauthor={Stephen Nelson}, pdfsubject={Gravity Field Dynamics (GFD)}, pdfkeywords={GFD, closure surface, field object, Jacobi surface, Hill problem, separatrix, vortex, Oort cloud, wide binaries, proof}}

\begin{document}

\articletype{Supplementary Material}

\title{Where Closure Sits: the Field's Stagnation Surface and the Jacobi Surface}

\author{Stephen Nelson}

\affil{Field Dynamics Research}

\ead{stephen@fielddynamics.org}

\begin{abstract}
Every field object closes, locally, at every level: an asteroid is a
pointer in the solar system's accounting, the solar system in the
galaxy's~\cite{NelsonMS}. This note asks where an object's closure
surface sits. In the frame turning with its orbit, the object's own
field, the host's tide and the orbit's turning make one effective
potential. Its field vanishes at exactly two points, the stagnation
points on the host's radius at the Jacobi radius, and the level through
them defines a surface from the field alone. Because the Coriolis force
does no work, the same surface is the one that matter below the
critical Jacobi energy can never leave: it stays in its cell. Matter
above it can snap to a neighbouring cell, and only through the necks
that open at the two stagnation points. So the surface the field defines
and the surface matter cannot leave are one surface. Without the host's
tide there are no stagnation points and no boundary: the surrounding
flow anchors the cell. A vortex in a strain
has the same structure: its closed streamlines are bounded by the
separatrix through its stagnation points. Under ownership the host's
tide is the mountain's and the object's own field is Newtonian, so the
rule reproduces itself. The surface is not a sphere: for the Sun it
reaches $1.44$~pc along the Galactic radius, $0.96$~pc along the orbit
and $0.73$~pc out of the plane. It holds Kepler orbits with $a < r_J/3
\approx 10^5$~AU, the outer Oort cloud's scale; it does not hold
$\alpha$~Centauri, though the two surfaces overlap; it predicts an edge
in the speeds of bound wide binaries at $\sqrt{2 - 3s/r_J}$ times the
circular speed. What remains a premise narrows to one geometric
statement: registration happens at the field's stagnation surface.
\end{abstract}

\keywords{GFD, closure surface, field object, Jacobi surface, Hill problem, separatrix, vortex, Oort cloud, wide binaries, proof}

% ----------------------------------------------------------------------

\section*{Setting}

GFD's Paper~I describes how a field object meets the NET Field: as a
whole, at its closure boundary, where settlement records one registered
state~\cite{NelsonI}; closure boundaries overlap and neighbours read
each other through them~\cite{NelsonII}. The mountain-scale proof writes
this into the action: each object's own field $d_k$ is removed from the
law's term, the field equation splits exactly, and nested closures
telescope, so each object is a pointer above its own level. Every object
closes; a drain decides only what its interior contains. It takes the
closure surface at the Jacobi radius as a premise~\cite{NelsonMS}.

Here the object is drainless (a star system, a binary, a cluster) and
moves on a near-circular orbit in its host. In the frame turning with
the orbit at $\Omega$, with $x$ along the host's radius, $y$ along the
orbit and $z$ out of the plane, and the object's size small against the
host's scale (Hill's limit),
\begin{equation}
\Phi_{\rm eff} = -\frac{GM}{r} - \frac12 T x^2 + \frac12\nu^2z^2,\qquad
T = 4\Omega^2 - \kappa^2,
\label{eq:phi}
\end{equation}
$\kappa$ the epicyclic and $\nu$ the vertical frequency. The members
move by $\ddot{\mathbf r} = -\nabla\Phi_{\rm eff} - 2\boldsymbol\Omega
\times\dot{\mathbf r}$~\cite{BT2008}.

Every statement is checked by script (16 checks); the algebra is
machine-checked in Lean~4 (8 theorems), importing the mountain-scale
formalisation~\cite{NelsonScript}.

% ----------------------------------------------------------------------

\section*{Theorem 1 (the field's stagnation surface)}

\emph{The combined field $-\nabla\Phi_{\rm eff}$ vanishes at exactly two
points, $x = \pm r_J$ on the host's radius, with
\[
r_J^3 = \frac{GM}{4\Omega^2 - \kappa^2},
\]
the Jacobi radius. They are saddles, with curvatures $(-3T, T, T +
\nu^2)$, and the level through them is $E_c = -\tfrac32GM/r_J$. Define
$\Sigma_{\rm field}$ as the boundary of the part of $\{\Phi_{\rm eff} <
E_c\}$ that contains the object. It is defined by the field alone, with
no orbits in it.}

\begin{proof}
The $y$ and $z$ components are $-GMy/r^3$ and $-(\nu^2 + GM/r^3)z$,
zero only on $y = z = 0$; there, $Tx = GM/x^2$. Lean
\texttt{stagnation\_radius}, \texttt{critical\_level}; script A1--A2.
\end{proof}

% ----------------------------------------------------------------------

\section*{Theorem 2 (the surface matter cannot leave)}

\emph{(a) The Coriolis force does no work, so the Jacobi energy $E_J =
\tfrac12v^2 + \Phi_{\rm eff}$ is conserved. Matter with $E_J < E_c$
inside $\Sigma_{\rm field}$ keeps $\Phi_{\rm eff} \le E_J < E_c$ and
never reaches the stagnation level: it stays in its cell and cannot snap
to a neighbouring one. The surface it cannot leave,
$\Sigma_{\rm Jacobi}$, is $\Sigma_{\rm field}$. (b) Above $E_c$ the
level opens in two necks around the stagnation points, and matter can
snap to a neighbouring cell only through them. The condition is
sufficient, not necessary: some matter above $E_c$ stays a long time,
retrograde orbits especially. (c) The cell is anchored by its
surroundings: as the host's tide $T \to 0$, $r_J = (GM/T)^{1/3} \to
\infty$, and with no surrounding flow there are no stagnation points
and no boundary.}

\begin{proof}
(a) Lean \texttt{coriolis\_no\_work}, \texttt{held\_in\_cell}. Script
B2: 40 orbits below $E_c$ over forty host orbits each stay in the cell
(farthest $0.74\,r_J$; $E_J$ conserved to $10^{-5}$). (b) Script B3:
just above $E_c$, $70\%$ of prograde and $3\%$ of retrograde orbits
snap out in that time, as known for Hill's problem~\cite{Henon1970},
and every one that left passed $1.05\,r_J$ within $30^\circ$ of the
host's radius, through the necks. (c) Lean \texttt{anchored}; script
A3.
\end{proof}

So the reviewers' question, whether registration and orbital binding
select the same surface, has a direct answer: they are the same level
set of the same potential, one read from the field, the other from the
matter. Part~(c) says the boundary is not the object's alone: it exists
only because of the field around it. In Paper~VII's picture the
unbound envelopes, their TetB chambers empty, form the mesh, and bound
envelopes, the matter, move through it~\cite{NelsonVII}; the mesh,
vorticing into TetB, is the surrounding flow that anchors a bound cell.
The proof's content is the tide's role; that this tide is the mesh's
flow is the framework's reading.

% ----------------------------------------------------------------------

\section*{Theorem 3 (the vortex)}

\emph{A point vortex of circulation $\Gamma$ in a pure strain of rate
$S$ has the streamfunction $\psi = -(\Gamma/2\pi)\ln r + Sxy$. The flow
stagnates at two points, $a^2 = \Gamma/4\pi S$ on the diagonal, which
are saddles of $\psi$. Fluid moves along $\psi = $ const, so the closed
streamlines around the vortex are bounded by the separatrix through the
stagnation points, and a parcel's fate is fixed by its $\psi$.}

\begin{proof}
Script C1: the Hessian determinant at the stagnation point is $-4$;
predicting 55 parcels' fates from $\psi$ alone agrees with integrating
them in every case (14 inside, 41 outside).
\end{proof}

This is the same structure one dimension lower: $\ln r$ for $1/r$, the
strain for the tide, $\psi$ for $\Phi_{\rm eff}$, and the separatrix for
$\Sigma_{\rm field}$. An eddy in a stream carries the fluid inside its
separatrix as one unit~\cite{Batchelor1956}; a bound object carries what
is inside its stagnation surface. The difference is the order of the
motion: fluid parcels follow $\psi$ exactly, while orbits are only
bounded by $\Phi_{\rm eff}$, which is why Theorem~2 is an inequality.

% ----------------------------------------------------------------------

\section*{Theorem 4 (shape and size)}

\emph{$\Sigma_{\rm field}$ is not a sphere. It reaches $r_J$ along the
host's radius (at the saddles), $\tfrac23r_J$ along the orbit (where
$-GM/y = E_c$) and, for the Sun's vertical frequency ($\nu^2 = 3.8\,T$,
local density $0.1\,M_\odot$\,pc$^{-3}$), $0.50\,r_J$ out of the plane.
Its volume is that of a sphere of radius $0.64\,r_J$.}

\begin{proof}
Lean \texttt{orbit\_extent}; script D1--D2.
\end{proof}

The mountain-scale proof's sphere of radius $r_J$ is therefore an outer
bound. Its results hold for any closed surface (balayage works onto any
regular boundary), and its Theorem~9 bounds the correction to the
Cassini quadrupole from a surface that is not a sphere~\cite{NelsonMS}.

% ----------------------------------------------------------------------

\section*{Theorem 5 (the rule reproduces itself)}

\emph{Expanding the host's potential in the turning frame about the
orbit, the radial curvature is $\Phi_h'' - \Omega^2 = -(4\Omega^2 -
\kappa^2)$, with $\kappa^2 = R\,d\Omega^2/dR + 4\Omega^2$; for a flat
rotation curve, the deep law's mountain, $T = 2\Omega^2$. Under
ownership the host's field is the mountain and the object's own field is
Newtonian. So $\Phi_{\rm eff}$ is exactly the potential the members move
in, and the surface they cannot cross is the one the field defines:
given ownership at $\Sigma_{\rm field}$, the dynamics hold every member
inside it.}

\begin{proof}
Script E1--E2; the split is the mountain-scale proof's Theorem~1.
\end{proof}

This is a fixed point, not a derivation from nothing: ownership at the
stagnation surface gives Newtonian internal motion, and Newtonian
internal motion respects exactly that surface. A different law inside
would move the surface, which is what makes the tests below
discriminate.

% ----------------------------------------------------------------------

\section*{Proposition 6 (tests)}

\emph{(a) The planets: for Neptune $E/E_c = r_J/3a = 3300$. (b) The
Oort cloud: a Kepler orbit is held for good when $-GM/2a < E_c$, i.e.
$a < r_J/3 = 99{,}000$~AU ($0.48$~pc), the scale usually quoted for the
outer Oort cloud; beyond it the Galactic tide strips comets over
time~\cite{Heisler1986}. (c) $\alpha$~Centauri: its surface and the
Sun's overlap, but its energy relative to the Sun, $+5\times10^8$~J/kg,
is far above the pair's $E_c = -9\times10^3$~J/kg. Neither belongs to
the other; each is its own cell, and they read each other's fields
through the overlap. (d) Wide binaries: a pair of separation $s$ is one
cell when $v/\sqrt{GM/s} < \sqrt{2 - 3s/r_J}$, an edge at $1.41$,
$1.40$, $1.38$, $1.35$, $1.28$ for $1.5\,M_\odot$ pairs at $s = 0.005$,
$0.02$, $0.05$, $0.1$, $0.2$~pc. Inside, the pair is drainless and its
relative motion Newtonian with $G_{\rm N}$. (e) Palomar~14: its surface
reaches $157$~pc along the Galactic radius and $104$~pc along its orbit,
around a $26$~pc body: one cell.}

\begin{proof}
Lean \texttt{held\_for\_good}, \texttt{binary\_edge}; script F1--F5.
\end{proof}

The wide-binary edge is set before any velocity is inspected: it uses
only the masses, the separation and the Galactic tide. GFD predicts that
the three-dimensional speed ratio of bound pairs stops there; a theory
that strengthens gravity inside the pair, such as MOND applied pointwise
with the external-field effect, moves the edge up. Gaia analyses are
divided on whether such a strengthening is seen~\cite{Chae2023,
Banik2024}; this proof does not adjudicate them, and observed speeds are
projected, which spreads the edge.

% ----------------------------------------------------------------------

\section*{Scope}

Proved here: the field's stagnation surface and the surface matter
cannot leave are one surface; matter snaps to a neighbouring cell only
through the necks at the stagnation points; the cell is anchored by its
surroundings; for a drainless object on a near-circular orbit
in Hill's limit; its shape and size; its vortex counterpart; that
ownership at it reproduces itself; and five tests.

Not proved here:
\begin{enumerate}
\item \emph{That registration happens at the stagnation surface.} This
is now the single geometric premise: the closure surface is the field's
separatrix. It is not derived from the dual-tetrad dynamics.
\item \emph{The snap itself}: what happens to the registered state on
both sides when matter changes cells, treated as a discrete event.
\item \emph{Eccentric orbits and a changing host}: the surface then
breathes, and ownership changes in time (stripping, capture).
\item \emph{Beyond Hill's limit}: when the object is not small against
its host's scale.
\item \emph{Orbits above $E_c$}: the energy condition is sufficient; how
long the others stay is measured, not bounded.
\item \emph{The drained level}: a galaxy's own closure surface in its
group or cluster, with the law acting inside.
\end{enumerate}

% ----------------------------------------------------------------------

\begin{thebibliography}{99}

\bibitem{NelsonMS}
S.~Nelson, \textit{The Mountain Scale in the SST: the Net-Field Rule
From the Action} (2026), \texttt{proofs/sst-mountain-scale/}.

\bibitem{NelsonI}
S.~Nelson, \textit{The Gravitational Constitutive Law}, Paper~I
(2026), \S\S2.2--2.3.

\bibitem{NelsonII}
S.~Nelson, \textit{A Covariant Theory Without Dark Matter}, Paper~II
(2026), \S6.

\bibitem{NelsonVII}
S.~Nelson, \textit{Dual Tetrad Gravity and the Field Origin},
Paper~VII (2026).

\bibitem{NelsonScript}
S.~Nelson, \texttt{verify\_closure\_surface.py} and
\texttt{lean/ClosureSurface.lean} (2026),
\texttt{proofs/closure-surface/}.

\bibitem{BT2008}
J.~Binney and S.~Tremaine, \textit{Galactic Dynamics}, 2nd ed.\
(Princeton University Press, 2008), \S8.3.

\bibitem{Henon1970}
M.~H\'enon, Numerical exploration of the restricted problem. VI. Hill's
case: non-periodic orbits, \textit{Astron.\ Astrophys.} \textbf{9}
(1970) 24.

\bibitem{Batchelor1956}
G.~K.~Batchelor, On steady laminar flow with closed streamlines at large
Reynolds number, \textit{J.\ Fluid Mech.} \textbf{1} (1956) 177.

\bibitem{Heisler1986}
J.~Heisler and S.~Tremaine, The influence of the Galactic tidal field
on the Oort comet cloud, \textit{Icarus} \textbf{65} (1986) 13.

\bibitem{Chae2023}
K.-H.~Chae, Breakdown of the Newton--Einstein standard gravity at low
acceleration in internal dynamics of wide binary stars,
\textit{Astrophys.\ J.} \textbf{952} (2023) 128.

\bibitem{Banik2024}
I.~Banik et al., Strong constraints on the gravitational law from Gaia
DR3 wide binaries, \textit{Mon.\ Not.\ R.\ Astron.\ Soc.} \textbf{527}
(2024) 4573.

\end{thebibliography}

\end{document}
